\section{Introducción: ¿por qué se necesitan modelos de lenguaje difusos}

Los modelos autoregresivos (Autoregressive, AR) han sido durante mucho tiempo el paradigma predominante en el campo de la generación de lenguaje. Desde la serie GPT hasta LLaMA, estos modelos generan texto palabra por palabra mediante un proceso estricto de izquierda a derecha (left-to-right), logrando un enorme éxito. Sin embargo, los modelos AR presentan una limitación fundamental de velocidad: generar $N$ tokens requiere $N$ pasos de propagación hacia adelante, lo que impide la generación verdadera en paralelo.

\begin{importantbox}{Ventaja central de los modelos de lenguaje difusos}
Los modelos de lenguaje difusos (Diffusion Language Models) se liberan de la restricción de izquierda a derecha y pueden generar texto en cualquier orden; lo más importante es que permiten \textbf{generar múltiples palabras en paralelo}. Esto les otorga una ventaja significativa en términos de velocidad de inferencia.
\end{importantbox}

Desde 2025, los modelos de lenguaje difusos han generado un interés generalizado en la industria. Google lanzó Gemini Diffusion, ByteDance introdujo Seed Diffusion e Inception Labs publicó Mercury Coder. Estos modelos no solo ofrecen un rendimiento excepcional, sino que superan ampliamente a los modelos AR tradicionales en velocidad de inferencia. Medios comerciales principales como Forbes y Fortune también han comenzado a reportar sobre esta transformación tecnológica.

El conferencista invitado de esta sesión, Subham Sahoo (萨哈姆·萨胡), es doctor en la Universidad Cornell; su investigación doctoral se centró en los modelos de lenguaje difusos. Es el primer autor de trabajos importantes como MDLM (Masked Diffusion Language Models) y SoLM (Sorting-based Language Models). Esta charla abordará a fondo las ideas centrales y los detalles técnicos de ambos artículos.

\subsection{Línea de tiempo del desarrollo de los modelos de lenguaje difusos}

\begin{knowledgebox}{De 2022 a 2025: hitos clave}
\begin{enumerate}
    \item \textbf{2022}: Se propuso por primera vez el modelo de lenguaje difuso. Estos modelos podían generar palabras en un orden aleatorio, pero la calidad del texto quedaba muy atrás respecto a los modelos AR, recibiendo una reacción fría por parte de la comunidad.
    \item \textbf{Principios de 2024}: Lou et al. demostraron que los modelos de lenguaje difusos combinados con arquitecturas Transformer modernas pueden igualar la calidad de los modelos AR a nivel de GPT-2, al mismo tiempo que son significativamente más rápidos.
    \item \textbf{Mediados de 2024}: Se publicó MDLM, que rederivó y simplificó el algoritmo de entrenamiento de difusión discreto, logrando un rendimiento comparable al de los modelos AR en conjuntos de datos de texto de pequeña escala.
    \item \textbf{Principios de 2025}: El artículo Diffusion Duality propuso un algoritmo de destilación que permite a los modelos generar secuencias de 1000 palabras con solo 10 pasos, manteniendo una calidad casi sin pérdida.
    \item \textbf{2025}: Se publicó SoLM, integrando AR y difusión discreta, con soporte para caché KV, lo que logró un aumento de orden de magnitud en la velocidad de inferencia.
\end{enumerate}
\end{knowledgebox}

\subsection{Resumen de la sección}

Los modelos de lenguaje difusos representan una transformación importante en el paradigma de generación de texto. Desde la propuesta conceptual en 2022 hasta el despliegue industrial en 2025, este campo ha experimentado un desarrollo rápido. El marco MDLM es actualmente la base de varios modelos de lenguaje difusos a nivel industrial (como Seed Diffusion de ByteDance).

\newpage
